\section{Reproducibility}
\label{app:repro}

\subsection{How the numbers and figures are regenerated}

Every training run saves its final metrics, a per-epoch training history, and
two checkpoints: the one with minimum training loss and the final one. A
single command-line entry point trains any configuration (model type, basis,
solver, sub-steps, learning rate, epoch budget and seed), and a separate
script re-scores any saved checkpoint against the ground truth it was
trained on. The extended runs of the budget validation were launched as
independent cloud jobs with the same entry point.

Everything reported after training is produced by an analysis pipeline that
retrains nothing. Its first step re-integrates every noise-free
$\Delta t=0.1$ checkpoint and checks that the rollout reproduces the run's
stored extrapolation MSE (the worst relative deviation is
$6.2\times10^{-3}$, from cloud-trained checkpoints re-scored on a different
CPU). The remaining steps compute the order conditions, work-precision data,
stability regions, Newton equilibria, first-integral drift, cross-solver
matrix, field-error maps and training milestones reported in
Sections~\ref{sec:solver-results}--\ref{sec:numerics} and~\ref{sec:phase4}.
Table~\ref{tab:provenance} lists what each derived figure and table is
computed from.

\begin{table}[t]
\centering
\footnotesize
\setlength{\tabcolsep}{4pt}
\renewcommand{\arraystretch}{1.08}
\caption{What each figure and table produced after training is computed from.}
\label{tab:provenance}
\begin{tabularx}{\columnwidth}{@{}L{4.1cm} Y@{}}
\toprule
\hdrrow \thh{Figure or table} & \thh{Computed from}\\
\midrule
Tables~\ref{tab:solver-ablation}, \ref{tab:basis-ablation}, \ref{tab:noise-dt}, \ref{tab:kan-vs-mlp}, \ref{tab:cost}, \ref{tab:phase4} & saved final metrics of each run\\
Table~\ref{tab:order-conditions} & the implemented solver tableaux\\
Figure~\ref{fig:work-precision}, truncation column of Table~\ref{tab:solver-ablation} & the true Lotka-Volterra field\\
Figure~\ref{fig:stability} & solver tableaux and the Tsit5 checkpoint\\
Figure~\ref{fig:basis-catalog}, $\kappa_2$ in Table~\ref{tab:basis-ablation} & the implemented basis functions\\
Table~\ref{tab:equilibria}, Figure~\ref{fig:first-integral} & trained checkpoints\\
Figure~\ref{fig:limit-cycle} & trained checkpoints\\
Figure~\ref{fig:cross-solver}, modified-equation analysis & checkpoints of the solver sweep\\
Figure~\ref{fig:field-error} & trained checkpoints\\
Figure~\ref{fig:edges} & RBF and B-spline checkpoints\\
Table~\ref{tab:milestones}, budget-validation crossover & per-epoch training histories\\
Figure~\ref{fig:pareto} & saved final metrics of each run\\
Figure~\ref{fig:noise-dt} (right) & runs of the noise sweep\\
Figure~\ref{fig:error-vs-time} & trained checkpoints\\
Diagrams (Figures~\ref{fig:kan-ode-architecture}, \ref{fig:kdense}, \ref{fig:data-split}, \ref{fig:rk-stages}, \ref{fig:locality}) & drawn from the implemented coefficients and shapes\\
\bottomrule
\end{tabularx}
\end{table}

\subsection{The integrator implementation}

The fixed-step Tsit5 step and the trajectory integrator, with docstrings
shortened. Every other method follows the same pattern with its own
tableau.

\begin{pycode}[Fixed-step Tsit5 step]
def step_tsit5(func, t, y, dt):
    # 6 of 7 stages: the 7th (FSAL) stage has TSIT5_B[6] == 0.0
    k = []
    k.append(func(t, y))
    for stage_idx in range(1, 6):
        t_stage = t + TSIT5_C[stage_idx] * dt
        y_stage = y
        for j in range(stage_idx):
            y_stage = y_stage + dt * TSIT5_A[stage_idx][j] * k[j]
        k.append(func(t_stage, y_stage))
    y_next = y
    for j in range(len(k)):
        if TSIT5_B[j] != 0.0:
            y_next = y_next + dt * TSIT5_B[j] * k[j]
    return y_next
\end{pycode}

\begin{pycode}[Trajectory integrator]
def odeint(func, y0, t, method="tsit5", substeps=1):
    step_fn = STEP_SOLVERS[method.lower()]
    trajectory = [y0]
    curr_y = y0
    for i in range(len(t) - 1):
        dt_sub = (t[i + 1] - t[i]) / substeps     # h = dt / substeps
        for s in range(substeps):
            curr_t = t[i] + s * dt_sub
            curr_y = step_fn(func, curr_t, curr_y, dt_sub)
        trajectory.append(curr_y)
    return torch.stack(trajectory, dim=0)
\end{pycode}

Because the loop is ordinary PyTorch, reverse-mode differentiation of the loss
records every stage evaluation; this is the ``direct autograd'' gradient
method of Section~\ref{sec:track-a}, and the reason training cost is
proportional to the number of function evaluations.

\subsection{Order-condition check}

The verification of Table~\ref{tab:order-conditions} needs only the
tableau. Rooted trees are represented as sorted tuples of their subtrees;
for a tree $\tau=[\tau_1,\dots,\tau_m]$ the elementary weight vector is
$\mathbf{g}(\tau)=\prod_k A\,\mathbf{g}(\tau_k)$ (elementwise, with
$\mathbf{g}(\bullet)=\mathbf{1}$), the density is
$\gamma(\tau)=|\tau|\prod_k\gamma(\tau_k)$, and the condition is
$\mathbf{b}^\top\mathbf{g}(\tau)=1/\gamma(\tau)$.

\begin{pycode}[Order-condition check]
def elem_weight(t, A):
    g = np.ones(A.shape[0])
    for child in t:
        g = g * (A @ elem_weight(child, A))
    return g

def gamma_of(t):
    return order_of(t) * math.prod(gamma_of(c) for c in t)

for q in range(1, 7):                  # 1, 1, 2, 4, 9, 20 trees
    devs = [b @ elem_weight(t, A) - 1.0 / gamma_of(t) for t in trees[q]]
\end{pycode}
